\section{Homework Week 1}

\begin{exercise}
    Let $R$ be a ring, and let $M$ be an $R$-module. Suppose that $S$ is a ring and
    \[
        \theta : S \to R
    \]
    is a ring homomorphism satisfying $\theta(1_S)=1_R$. Let $S$ act on $M$ by
    \[
        sx=\theta(s)x,
    \]
    where $x\in M$ and $s\in S$. Prove that $M$ is an $S$-module.
\end{exercise}

\begin{proof}
    This action is well-defined because $\forall s \in S$,$\theta(s) \in R$ and $\theta(s)x \in M$ for all $x \in M$. We need to verify the module axioms, for all $s, t \in S$ and $x, y \in M$:
    \begin{itemize}
        \item $s(x+y) = \theta(s)(x+y) = \theta(s)x + \theta(s)y = sx + sy$ \item $(s+t)x = \theta(s+t)x = \theta(s)x + \theta(t)x = sx + tx$ \item $(st)x=\theta(st)x=\theta(s)\theta(t)x=\theta(s)(\theta(t)x)=s(tx)$.
        \item $1_S x = \theta(1_S)x = 1_R x = x$ \end{itemize}
    Hence $M$ is an $S$-module.
\end{proof}

\begin{exercise}
    Let $R$ be a ring, and let $M$ and $N$ be two right $R$-modules. Let $\operatorname{Hom}_{\mathbb{Z}}(M,N)$ be the set of group homomorphisms from $M$ to $N$(as abelian groups). The action of $R$ on $\operatorname{Hom}_{\mathbb{Z}}(M,N)$ is given as follows: for any $a\in R$ and $\varphi\in \operatorname{Hom}_{\mathbb{Z}}(M,N)$, define
    \[
        (\varphi a)(x)=\varphi(xa), \qquad \forall x\in M.
    \]
    Prove that $\operatorname{Hom}_{\mathbb{Z}}(M,N)$ is a left $R$-module.
\end{exercise}

\begin{proof}
    This action is well defined: for every $r \in R$ and every $\varphi \in \operatorname{Hom}_{\mathbb{Z}}(M,N)$, the map
    \[
        \varphi r\colon M \to N, \qquad (\varphi r)(x)=\varphi(xr),
    \]
    is again a group homomorphism.

    We now verify the module axioms. Let $r_1, r_2 \in R$, let $\varphi,\psi \in \operatorname{Hom}_{\mathbb{Z}}(M,N)$, and let $x \in M$. Then
    \begin{align*}
        ((\varphi+\psi)r_1)(x)
         & = (\varphi+\psi)(xr_1)               \\
         & = \varphi(xr_1)+\psi(xr_1)           \\
         & = (\varphi r_1)(x)+(\psi r_1)(x),    \\
        (\varphi(r_1+r_2))(x)
         & = \varphi(x(r_1+r_2))                \\
         & = \varphi(xr_1+xr_2)                 \\
         & = \varphi(xr_1)+\varphi(xr_2)        \\
         & = (\varphi r_1)(x)+(\varphi r_2)(x), \\
        ((\varphi r_1)r_2)(x)
         & = (\varphi r_1)(xr_2)                \\
         & = \varphi((xr_2)r_1)                 \\
         & = \varphi(x(r_2r_1))                 \\
         & = (\varphi(r_2r_1))(x),              \\
        (\varphi 1_R)(x)
         & = \varphi(x1_R)                      \\
         & = \varphi(x).
    \end{align*}
    Hence $\operatorname{Hom}_{\mathbb{Z}}(M,N)$ is a left $R$-module.
\end{proof}

\begin{exercise}
    Let $R$ be a ring, and let $M$ be an $R$-module. Denote
    \[
        \operatorname{Ann}_R M=\{r\in R \mid rm=0,\ \forall m\in M\}.
    \]
    We call $M$ a faithful $R$-module if $\operatorname{Ann}_R M=\{0\}$. Define the action
    of the quotient ring $R/\operatorname{Ann}_R M$ on $M$ by
    \[
        \overline{r}\,x=rx,
    \]
    where $x\in M$,$r\in R$, and $\overline{r}=r+\operatorname{Ann}_R M$. Prove that $M$ is a faithful $R/\operatorname{Ann}_R M$-module.
\end{exercise}

\begin{proof}
    First we show that the action is well-defined. Suppose
    \[
        \bar r=\bar r' \in R/\operatorname{Ann}_R M.
    \]
    Then $r-r'\in \operatorname{Ann}_R M$, so for every $x\in M$,
    \[
        (r-r')x=0.
    \]
    Hence $rx=r'x$, so $\bar r\,x$ does not depend on the choice of representative.

    The module axioms follow immediately from those of the $R$-module structure on $M$.

    Now we prove faithfulness. Suppose $\bar r\in R/\operatorname{Ann}_R M$ acts trivially on $M$, that is,
    \[
        \bar r\,x=0 \qquad \forall x\in M.
    \]
    By definition, this means
    \[
        rx=0 \qquad \forall x\in M.
    \]
    Therefore $r\in \operatorname{Ann}_R M$, so $\bar r=\bar 0$ in the quotient ring. Hence
    \[
        \operatorname{Ann}_{R/\operatorname{Ann}_R M}(M)=\{\bar 0\},
    \]
    and $M$ is faithful as an $R/\operatorname{Ann}_R M$-module.
\end{proof}

\begin{exercise}
    Let $\varphi : M \to M$ be a homomorphism of $R$-modules such that
    \[
        \varphi\varphi=\varphi.
    \]
    Prove that
    \[
        M=\ker\varphi \oplus \operatorname{Im}\varphi.
    \]
\end{exercise}

\begin{proof}
    Let $x \in M$. Then
    \[
        x = \varphi(x) + (x - \varphi(x)),
    \]
    where $\varphi(x) \in \operatorname{Im}\varphi$ and $x - \varphi(x) \in \ker\varphi$. Hence
    \[
        M = \operatorname{Im}\varphi + \ker\varphi.
    \]
    It remains to show that
    \[
        \ker\varphi \cap \operatorname{Im}\varphi = \{0\}.
    \]
    Suppose $x \in \ker\varphi \cap \operatorname{Im}\varphi$. Then $x = \varphi(y)$ for some $y \in M$ and $\varphi(x) = 0$. Thus
    \[
        0 = \varphi(x) = \varphi(\varphi(y)) = \varphi(y)=x,
    \]
    so $x = 0$. Hence $\ker\varphi \cap \operatorname{Im}\varphi = \{0\}$.
\end{proof}

\begin{exercise}
    Show that any $R$-module is a homomorphism image of some free $R$-module.
\end{exercise}

\begin{proof}
    Let $M$ be a left $R$-module. Consider the free $R$-module
    \[
        F=\bigoplus_{x\in M} Re_x
    \]
    with basis $\{e_x\}_{x\in M}$ indexed by the underlying set of $M$.
    Define an $R$-module homomorphism
    \[
        \pi:F\to M
    \]
    by
    \[
        \pi(e_x)=x \qquad (x\in M),
    \]
    and extend $R$-linearly. Thus
    \[
        \pi\Bigl(\sum_{i=1}^n r_i e_{x_i}\Bigr)=\sum_{i=1}^n r_i x_i.
    \]
    This is well-defined and $R$-linear.

    For any $m\in M$, we have
    \[
        \pi(e_m)=m,
    \]
    so $\pi$ is surjective. Therefore $M$ is a homomorphic image of the free module $F$.
    Equivalently,
    \[
        M\cong F/\ker\pi.
    \]
\end{proof}

\begin{exercise}
    Let $\mathbb{Q}$ be the field of rational numbers. Is the additive group $(\mathbb{Q},+)$ a free $\mathbb{Z}$-module?
\end{exercise}

\begin{proof}
    Claim: any two non-zero element of $\mathbb{Q}$ are linearly dependent over $\mathbb{Z}$. Hence if $\mathbb{Q}$ is a free $\mathbb{Z}$-module, then it must have a
    basis of cardinality 1.

    Proof of claim: let $0 \neq x = \frac{a}{b}$ and $0 \neq y =\frac{c}{d}$ be two
    elements of $\mathbb{Q}$. Then $dcx - aby=0$. If $a=0$ or $c=0$, then $x=0$ or $y=0$,
    which contradicts the assumption. Hence $a, c \neq 0$ which means $dc,ab$ are two
    non-zero integer coefficients. Hence $x$ and $y$ are linearly dependent over $\mathbb{Z}$. And if $\mathbb{Q}$ is a free $\mathbb{Z}$-module, then it must have a
    basis of cardinality 1.

    But if $\mathbb{Q}$ is a cyclic $\mathbb{Z}$-module, then it has only one generator.
    Suppose it is $z=\frac{e}{f}$. Then $\mathbb{Q}=\{n \frac{e}{f} \mid n \in \mathbb{Z}\}$.
    But this is not possible because we can easily pick $\frac{1}{2f} \notin \mathbb{Q}$,
    which is absurd.

    Therefore $\mathbb{Q}$ is not a free $\mathbb{Z}$-module.
\end{proof}

